\scanpage{52}
(ii) Par une construction reposant \unsure{ess\textsuperscript{t}|ess\textsuperscript{l}} sur le même principe, mais seulement un peu plus sophistiquée techniquement, (cf appendice), on trouve que si $f$ est lisse \unsure{qu. cpct|qu. cpt} et qu. sép., alors
\[
\mathbb R f_*(\Omega^{\cdot}_{X/S})
\]
est muni d'une \uline{connexion} \uline{absolue canonique} [au sens des catégories dérivées]. Cela ne signifie pas tt. à fait que les $\mathbb R^i f_*(\Omega^{\cdot}_{X/S})$ sont munis individuell\textsuperscript{t} des connexions absolues ; c'est pourtant vrai si $S$ de car. $0$, et $f$ propre et lisse (car alors la formation des $R^i f_*(\Omega^{\cdot}_{X/S})$ eux-mêmes commute à l'extension de la base ....) -- ou encore, bien sûr, dans le cas lisse affine envisagé en (i) .... La vérification (pour $f$ propre, lisse) que cette connexion est compatible avec la description transcendante, est en fait triviale.

\uline{Remarques} On \unsure{pourrait|peut} espérer que ces connexions canoniques proviennent en fait de \uline{stratifications} canoniques, de façon que l'on aurait une
% archive sheet number: 51

